\documentclass[10pt]{article}
\usepackage[T1]{fontenc}
\usepackage[utf8]{inputenc}
\usepackage[french]{babel}
\usepackage{amsmath,amssymb,mathrsfs,mathpazo}
\usepackage{geometry}
\geometry{a4paper,margin=2cm}
\usepackage{graphicx,array,tabularx,tabularray,colortbl,multirow,multicol,booktabs,enumitem,pifont,ulem,fancybox,diagbox,cancel}
\usepackage{pstricks,pstricks-add,pst-plot,pst-tree,pst-node,pst-text,pst-func,pst-math,pst-bspline,pst-3dplot,pst-eucl,pst-solides3d}
\let\pscalculate\undefined
\usepackage{pst-bezier}
\usepackage{tikz,pgfplots,tkz-tab}
\usetikzlibrary{arrows,arrows.meta,patterns,calc,positioning,trees,decorations.markings,shapes.geometric,shapes.misc,shapes.arrows,intersections,decorations.pathmorphing,decorations.pathreplacing,backgrounds,shadows,babel}
\pgfplotsset{compat=1.18}
\usepgfplotslibrary{fillbetween}
\usepackage[np]{numprint}
\usepackage{listings,makecell,pas-tableur,twoopt}
\let\barre\undefined
\usepackage[pstricks]{tablvar}
\definecolor{darkgreen}{RGB}{0,100,0}
\providecommand{\red}{\color{red}}
\providecommand{\blue}{\color{blue}}
\providecommand{\green}{\color{green}}
\providecommand{\black}{\color{black}}
\providecommand{\white}{\color{white}}
\providecommand{\gray}{\color{gray}}
\providecommand{\cyan}{\color{cyan}}
\providecommand{\magenta}{\color{magenta}}
\providecommand{\yellow}{\color{yellow}}
\providecommand{\violet}{\color{violet}}
\providecommand{\orange}{\color{orange}}
\providecommand{\euro}{\texteuro}
\providecommand{\EUR}{\texteuro}
\providecommand{\ieme}{\textsuperscript{e}}
\providecommand{\degres}{\ensuremath{^\circ}}
\providecommand{\degree}{\ensuremath{^\circ}}
\providecommand{\vv}[1]{\overrightarrow{#1}}
\providecommand{\Coord}[1]{#1}
\providecommand{\decosix}{}
\providecommand{\starredbullet}{\ensuremath{\bullet}}
\pagestyle{empty}
\setlength{\parindent}{0pt}
\usepackage{adjustbox}\DeclareUnicodeCharacter{2236}{:}\DeclareUnicodeCharacter{2190}{\ensuremath{\leftarrow}}
\providecommand{\R}{\mathbb{R}}
\providecommand{\N}{\mathbb{N}}
\providecommand{\Z}{\mathbb{Z}}
\providecommand{\Q}{\mathbb{Q}}
\providecommand{\D}{\mathbb{D}}
\providecommand{\C}{\mathbb{C}}
\providecommand{\e}{{\mathrm{e}}}
\providecommand{\vect}[1]{\overrightarrow{#1}}
\providecommand{\vectt}[1]{\overrightarrow{\mathrm{#1}}}
\providecommand{\vvt}[1]{\overrightarrow{\mathrm{#1}}}
\providecommand{\np}[2][]{{#2~#1}}
\providecommand{\vv}[1]{{\overrightarrow{#1}}}
\renewcommand{\d}{{\mathrm{d}}}
\providecommand{\degres}{{^\circ}}
\providecommand{\degre}{{^\circ}}
\providecommand{\blue}{\color{blue}}
\providecommand{\red}{\color{red}}
\providecommand{\enskip}{\quad}
\providecommand{\llbracket}{\lbrack\!\lbrack}
\providecommand{\rrbracket}{\rbrack\!\rbrack}
\providecommand{\euro}{{\text{€}}}
\providecommand{\no}{{\text{n°}}}
\providecommand{\interval}[2]{{[#1;#2]}}
\providecommand{\intervFO}[2]{{[#1;#2[}}
\providecommand{\dotfill}{\ldots}

\providecommand{\ds}{\displaystyle}
\providecommand{\textding}[1]{#1}
\providecommand{\Oij}{\(\mathrm{O};\vec{\imath},\vec{\jmath}\)}
\providecommand{\Oijk}{\(\mathrm{O};\vec{\imath},\vec{\jmath},\vec{k}\)}
\providecommand{\Ouv}{\(\mathrm{O};\vec{u},\vec{v}\)}
\providecommand{\intervFF}[2]{\ensuremath{[#1;#2]}}
\renewcommand{\euro}{\ifmmode\text{\texteuro}\else\texteuro\fi}
\DeclareUnicodeCharacter{00B0}{\ifmmode^{\circ}\else\textdegree\fi}
\DeclareUnicodeCharacter{20AC}{\ifmmode\text{\texteuro}\else\texteuro\fi}
% Énoncés : Éducation nationale. Transcription et corrigés : APMEP, Denis Vergès et contributeurs. Découpage et classement : CoopMaths, Sébastien Lozano, Éric Elter et contributeurs. Adaptation de la présentation : Maths Explorer.
% CC BY-SA : https://coopmaths.fr/www/about/
% https://coopmaths.fr/alea/static/e3c/2020/tex/e3c_2020_00_sujet62_3_cor.tex
\begin{document}
\noindent\textit{APMEP / CoopMaths. CC BY-SA. Présentation adaptée par Maths Explorer.}\par\medskip
\subsection*{1.}

\begin{center}
\begin{tikzpicture}[thick, scale=1.5] 
\node (P_-1_0) at (-2,-1.5) {$\phantom{A}$};
\node (P_0_0) at (0,-0.5) {$N$};
\draw (P_-1_0) -- (P_0_0) node[midway, above] {$0{,}6$};
\node (P_1_0) at (2,-0) {$G$};
\draw (P_0_0) -- (P_1_0) node[midway, above] {$0{,}45$};
\node (P_1_1) at (2,-1) {$\overline{G}$};
\draw (P_0_0) -- (P_1_1) node[midway, below] {$0{,}55$};
\node (P_0_2) at (0,-2.5) {$\overline{N}$};
\draw (P_-1_0) -- (P_0_2) node[midway, below] {$0{,}4$};
\node (P_1_2) at (2,-2) {$G$};
\draw (P_0_2) -- (P_1_2) node[midway, above] {$0{,}7$};
\node (P_1_3) at (2,-3) {$\overline{G}$};
\draw (P_0_2) -- (P_1_3) node[midway, below] {$0{,}3$};
\end{tikzpicture}
\end{center}

\subsection*{2.}

On a :
\[
p(N \cap G) = p(N) \times p_N(G) = 0{,}6 \times 0{,}45 = 0{,}27.
\]

\subsection*{3.}

On a de même :
\[
p(\overline{N} \cap G) = p(\overline{N}) \times p_{\overline{N}}(G) = 0{,}4 \times 0{,}7 = 0{,}28.
\]
D'après la loi des probabilités totales :
\[
p(G) = p(N \cap G) + p(\overline{N} \cap G) = 0{,}27 + 0{,}28 = 0{,}55.
\]

\subsection*{4.}

Il faut trouver :
\begin{align*}
p_{\overline{G}}(N) &= \dfrac{p(\overline{G} \cap N)}{p(\overline{G})} \\
&= \dfrac{p(N \cap \overline{G})}{1 - p(G)} \\
&= \dfrac{0{,}6 \times 0{,}55}{1 - 0{,}55} \\
&= \dfrac{0{,}33}{0{,}45} \approx 0{,}733.
\end{align*}

\subsection*{5.}

On peut dresser un arbre représentant les résultats des deux combats :

\begin{center}
\begin{tikzpicture}[thick, scale=1.5] 
\node (P_-1_0) at (-2,-1.5) {$\phantom{A}$};
\node (P_0_0) at (0,-0.5) {$G$};
\draw (P_-1_0) -- (P_0_0) node[midway, above] {$0{,}55$};
\node (P_1_0) at (2,-0) {$G$};
\draw (P_0_0) -- (P_1_0) node[midway, above] {$0{,}55$};
\node (P_1_1) at (2,-1) {$\overline{G}$};
\draw (P_0_0) -- (P_1_1) node[midway, below] {$0{,}45$};
\node (P_0_2) at (0,-2.5) {$\overline{G}$};
\draw (P_-1_0) -- (P_0_2) node[midway, below] {$0{,}45$};
\node (P_1_2) at (2,-2) {$G$};
\draw (P_0_2) -- (P_1_2) node[midway, above] {$0{,}55$};
\node (P_1_3) at (2,-3) {$\overline{G}$};
\draw (P_0_2) -- (P_1_3) node[midway, below] {$0{,}45$};
\end{tikzpicture}
\end{center}

On a :

\( p(G \cap G) = 0{,}55 \times 0{,}55 = 0{,}3025 \) ;

\( p\left(\overline{G} \cap \overline{G}\right) = 0{,}45 \times 0{,}45 = 0{,}2025 \).

Il reste donc pour la probabilité d'une victoire et une défaite (ou inversement) : 
\[
1 - (0{,}3025 + 0{,}2025) = 1 - 0{,}505 = 0{,}495.
\]
D'où le tableau de la loi de probabilité de \( X \) :

\[
\begin{array}{|c|c|c|c|}
\hline
\text{Valeur de } X : x_i & 0 & 1 & 2 \\
\hline
p(X = x_i) & 0{,}2025 & 0{,}495 & 0{,}3025 \\
\hline
\end{array}
\]

\end{document}
